\section{Related Work}
\label{sec:related-work}

% Full map with justification text: NARRATIVE.md \S5. Kept LIGHT per that
% plan — sharper distinctions (esp. vs. M2RNN, DeltaProduct) are deferred
% to forward-references from \S\ref{sec:mechanism}/\S\ref{sec:the-fix}
% where they are actually load-bearing.

\textbf{Descriptive rank measurement, the primary foil.} A recent wave of
work measures the effective rank of the fixed-size state maintained by
fast-weight linear-attention models on real, pretrained language, and
finds it bounded by sequence length \citep{nazari2026rank,
sun2026staterank}. Both papers use the same effective-rank instrument we
do, but leave $K$ uncontrolled, offer no necessity theorem, and run no
causal intervention. This paper is close to a mirror image of that
program on the same architecture family: controlled $K$, a provable lower
bound, and causal (both train-time and eval-time) intervention.

\textbf{Theoretical constructions.} The closest prior theoretical result
is \citet{nichani2025factual} (ICLR 2025 Spotlight), a rank-$m$ existence
proof under discrete argmax decoding, where a rank-1 state can recover
$\approx d$ associations. This is precisely what motivates pinning our
own readout to exact, continuous cosine recovery rather than argmax
(Section~\ref{sec:setup}) --- the argmax result would otherwise silently
defeat any necessity claim built on this architecture.
\citet{schlag2021linear} originates the ``linear-attention state as
associative memory'' framing this paper inherits, in the fast-weight-programmer
lineage. \citet{barnfield2026sharp} derive capacity thresholds for static,
rank-constrained associative memories with no recurrence and no training
dynamics; our contribution relative to that work is the training-time
force-rank ablation on a time-varying, streamed state, together with
held-out-hop composition. \citet{arora2024simple} (Based, Thm.\ F.1) and
the associated Zoology benchmark suite \citep{arora2023zoology} bound
state \emph{size} in a decoding-agnostic sense; under our pinned
linear-unbind readout, rank rather than size is the binding resource, and
our synthetic probe grammar is a close relative of Zoology's
associative-recall (MQAR) family.

\textbf{Architecture family.} We use \texttt{fla-org}'s production
DeltaNet / Gated-DeltaNet kernel directly. DeltaProduct
\citep{siems2025deltaproduct} extends the same lineage with a
per-step transition matrix $A_t = I - \beta_t k_t k_t^\top$ whose own
``rank'' is a different object from the state's rank $\mathrm{rank}(S_t)$
measured throughout this paper --- the single most likely point of
reviewer confusion, and worth stating explicitly rather than leaving
implicit. \citet{grazzi2025negative} study negative-eigenvalue state
tracking, whose eigenvalue sign/phase framing is relevant to this paper's
own $h{=}21$ depth-amplification discussion (Section~\ref{sec:phenomenon}).

\textbf{Superposition and distributed-representation capacity.}
\citet{elhage2022superposition}'s framing of superposition, and the
VSA/HRR capacity tradition of \citet{plate1995holographic} (Holographic
Reduced Representations) and \citet{frady2020resonator} (Resonator
Networks), all study a different notion of capacity --- SNR and bundling
under a favorable, typically sparse readout --- than the exact-recovery
requirement this paper's provable bound closes off. We scope the
contrast explicitly rather than let the terms collide.

\textbf{M\textsuperscript{2}RNN, the closest prior-art paper}
\citep{mishra2026m2rnn}. M\textsuperscript{2}RNN trains a matrix-valued
recurrent state on permutation-group ($S_3$) composition at sequence
length 128 and evaluates near-perfectly ($\geq99.5\%$ accuracy) at 512.
Leading with what their own results concede: their comparison table
already has a plain vector-state GRU matching the matrix model on this
exact task (also $\geq99.5\%$ at 512 --- rank is not shown load-bearing
there; the GRU only loses on general LM perplexity, a different axis),
and their Theorem~1 (the recurrence can represent all tasks representable
by non-linear vector-valued RNNs) is a representational claim, not a
necessity one --- it proves no lower bound against vector-valued models
and runs no causal rank intervention. This paper is the three-way mirror:
a provable lower bound (not just a representation result), a causal
intervention (not just an evaluation), and a controlled $K$ (not just a
fixed task). Both sub-claims above were independently re-verified against
the fetched paper text for this submission (see \texttt{refs.bib}).

\textbf{A credible rival causal account, engaged directly}
\citep{okpekpe2025revisiting}. That paper argues Transformers and SSMs
differ in optimization dynamics rather than expressive power --- MQAR-style
failures in Mamba and Hyena are gated by a narrow learning-rate window,
not a hard capacity ceiling, and disappear under a finer LR grid. A
hostile reading of this paper's own capacity cliff
(Section~\ref{sec:phenomenon}, $x_0{=}0.5455$) is that it, too, is an
optimization artifact rather than a structural one. We do not think it
is, and the reason is architectural, not rhetorical: our causal chain
runs through forced-rank construction and eval-time SVD truncation on an
\emph{already-trained} checkpoint, never through a learning-rate axis, and
the recovery ceiling lands exactly at $k{=}K$ regardless of how that
checkpoint was optimized --- a structural claim their account does not
make or predict.

\textbf{Frozen-QK, the nearest neighbor for the frozen-anchor-table
mechanism} \citep{dong2025frozenqk} (Sections~\ref{sec:mechanism},
\ref{sec:the-fix}). That work fully freezes an entire softmax
transformer's Q/K projections at initialization; induction heads still
form, but at a real cost ($\approx$10\% WikiText perplexity regression).
Our frozen anchor table differs on three axes: it is additive, not a
full-freeze --- a tuned-weight blend into an otherwise fully-trained key
path, not a projection replacement; its own gain is loss-neutral by
construction, since the same architecture with the table zeroed out is
the baseline; and it targets a bounded recurrent fast-weight state, not
unbounded-context softmax attention, so the question it answers (does a
fixed-size state retain more) has no direct analog in Frozen-QK's
setting.

\textbf{VLA, the geometric rival account for the capacity cliff}
\citep{pandey2026vla} (Section~\ref{sec:dose-response}). VLA's stabilized
architecture (a Kalman/variational update with a spectral-norm-1 write
direction) pushes a sharp MQAR-recall cliff outward, framing capacity as
tied to key linear independence and state-matrix rank --- a geometric
account in tension with our own direct coherence-injection result, which
finds $h4{=}1.0$ at every dose tested up to and beyond $d{=}64$'s realized
coherence band (10/10 cells). We state the tension rather than bury it:
VLA studies single-hop recall under an architectural stabilization, while
this paper studies relational binding and multi-hop composition under a
direct causal manipulation of coherence at matched load. It remains
possible both accounts describe the same underlying fact from different
angles --- a rank-deficient key population is, definitionally, also
capacity-constrained at fixed $d_{\mathrm{state}}$ --- and we do not
claim to have resolved which framing is more fundamental, only that our
own causal dose-response manipulation did not find coherence-as-such to
be the lever.

\textbf{Immediate lineage.} This paper transfers and extends our own
prior, unpublished internal work, released as part of this submission's
anonymized supplementary artifact rather than cited as a public
reference: a from-scratch, bespoke-attention-encoder version of the
identical rank-recruitment result, including the depth-amplification law
imported directly, and an earlier negative result on a discrete-decoding
readout that motivated the exact-continuous readout choice used
throughout this paper. What is new here is the transfer of that result to
a production architecture and, novel to this paper, the discovery and
fix of the exactness gap the synthetic construction never exposed.
